\section{Experiments}

In this section, we describe our experimental methodology for evaluating the SO(3) tube smoothing algorithm. We design experiments to validate both the supporting analysis and practical performance on both synthetic and real-world datasets.

\subsection{Experimental Setup}

\textbf{Evaluation Metrics:} We use four primary metrics to assess algorithm performance:

\begin{itemize}
\item \textit{Tube Constraint Satisfaction:} Maximum and average tube excess:
\begin{equation}
\text{tube\_excess}_{\text{max}} = \max_{i=0}^{N-1} \left( \left\| \log(R_i^T \hat{R}_i) \right\|_2 - \epsilon_i \right)
\end{equation}
\begin{equation}
\text{tube\_excess}_{\text{avg}} = \frac{1}{N} \sum_{i=0}^{N-1} \max\left( \left\| \log(R_i^T \hat{R}_i) \right\|_2 - \epsilon_i, 0 \right)
\end{equation}

\item \textit{Geodesic Error:} Angular distance to ground truth (when available):
\begin{equation}
\epsilon_{\text{geo}} = \frac{1}{N} \sum_{i=0}^{N-1} \left\| \log(R_{\text{true},i}^T \hat{R}_i) \right\|_2
\end{equation}

\item \textit{Smoothness Quality:} Root mean square of angular velocity and acceleration:
\begin{equation}
\omega_{\text{RMS}} = \sqrt{\frac{1}{N-1} \sum_{i=0}^{N-2} \left\| \phi_{i+1} - \phi_i \right\|_2^2}
\end{equation}
\begin{equation}
\alpha_{\text{RMS}} = \sqrt{\frac{1}{N-2} \sum_{i=1}^{N-2} \left\| \phi_{i-1} - 2\phi_i + \phi_{i+1} \right\|_2^2}
\end{equation}

\item \textit{Computational Performance:} Runtime, memory usage, and iteration counts for both outer and inner loops.
\end{itemize}

\textbf{Implementation Details:} All experiments use Python 3.11 with NumPy and SciPy for numerical computations. The ADMM solver uses cached sparse Cholesky factorization with fallback to preconditioned conjugate gradient. We implement the SOCP baseline using CVXPY with ECOS and SCS solvers~\cite{boyd_scp}.

\subsection{Synthetic Bounded-Noise Experiments}

We design synthetic experiments to validate theoretical properties and demonstrate practical tube compliance.

\textbf{Experiment 1: Feasibility Verification.} We generate smooth ground-truth trajectories with controlled noise levels and set tube radii $\epsilon_i$ such that:
\begin{equation}
\epsilon_i \geq \left\| \log(R_{\text{true},i}^T R_{\text{meas},i}) \right\|_2 + \delta_{\text{margin}}
\end{equation}
where $\delta_{\text{margin}} = 0.05$ rad provides safety margin. This ensures feasible solutions exist by construction.

We evaluate:
\begin{itemize}
\item \textbf{Tube compliance:} Maximum tube excess across all samples
\item \textbf{Algorithm convergence:} Number of outer iterations to reach $\|\delta\|_\infty < 10^{-6}$
\item \textbf{Warm-starting impact:} Comparison with and without warm-starting on consecutive smoothing problems
\end{itemize}

\textbf{Experiment 2: Outlier Detection.} We introduce controlled outliers by setting $\epsilon_i$ too small for specific samples (e.g., $\epsilon_k = 0.02$ rad while $\epsilon_i = 0.15$ rad for others). We evaluate the slack variable mechanism by monitoring:

\begin{itemize}
\item \textbf{Outlier identification:} Which samples activate slack ($s_i > 0$)
\item \textbf{Constraint preservation:} Whether non-outlier samples maintain feasibility
\item \textbf{Smoothness impact:} Angular velocity RMS with and without outlier handling
\end{itemize}

\textbf{Experiment 3: Scaling Analysis.} We evaluate computational performance across problem sizes:
\begin{equation}
N \in \{10^3, 10^4, 5 \times 10^4, 10^5\}
\end{equation}
For each size, we measure:
\begin{itemize}
\item \textbf{Runtime scaling:} Total time as function of $N$
\item \textbf{Iteration counts:} Outer iterations and average ADMM iterations
\item \textbf{Memory usage:} Peak memory consumption
\end{itemize}

\subsection{Real Sensor Data Experiments}

We evaluate on real sensor data to demonstrate practical applicability.

\textbf{Dataset: EuRoC MAV Machine Hall Sequences.} We use sequences from the EuRoC MAV dataset~\cite{euroc}, specifically:
\begin{itemize}
\item \textbf{MH\_01\_easy:} Machine hall traversal with moderate motion
\item \textbf{MH\_02\_easy:} Straight corridor navigation
\item \textbf{V1\_02\_medium:} Outdoor environment with varying lighting
\end{itemize}

Each sequence provides rotation estimates from visual-inertial odometry and ground truth from motion capture systems.

\textbf{Tube Radius Specification:} Unlike synthetic experiments where $\epsilon_i$ is set for feasibility, real data requires physically meaningful $\epsilon_i$ values. We derive $\epsilon_i$ from:

\begin{itemize}
\item \textbf{Sensor specifications:} IMU gyroscope bias and noise characteristics
\item \textbf{Calibration data:} Visual odometry reprojection error analysis
\item \textbf{Time-varying uncertainty:} Larger $\epsilon_i$ during fast motion or poor lighting
\end{itemize}

For example, for gyroscope data with bias $b_g$ and noise $\sigma_g$ over time step $\Delta t$:
\begin{equation}
\epsilon_i = 3\sqrt{b_g^2 + \sigma_g^2} \Delta t
\end{equation}
This provides physically meaningful bounds that reflect actual sensor capabilities rather than tuned parameters.

\textbf{Evaluation on Real Data:} We compare our method against:
\begin{itemize}
\item \textbf{Unconstrained baseline:} Same energy function without tube constraints
\item \textbf{SOCP baseline:} CVXPY with ECOS/SCS solvers (slower but for reference)
\item \textbf{Prior work:} Jia & Evans~\cite{jia_evans} constrained rotation smoothing
\end{itemize}

Metrics include geodesic error to ground truth, smoothness quality, and runtime.

\subsection{Baseline Methods}

We implement the following baselines for fair comparison:

\textbf{Method 1: Unconstrained Smoothing.} Solves the unconstrained optimization:
\begin{equation}
\min_\phi E(\phi) = \frac{\lambda}{2\tau}\sum_{i=0}^{N-2} \|\phi_{i+1} - \phi_i\|_2^2 + \frac{\mu}{2\tau^3}\sum_{i=1}^{N-2} \|\phi_{i-1} - 2\phi_i + \phi_{i+1}\|_2^2
\end{equation}
This provides optimal smoothness but no explicit tube-compliance guarantees.

\textbf{Method 2: SOCP Baseline.} Uses CVXPY to solve the same problem formulation as our method, but without exploiting block-banded structure. Solvers:
\begin{itemize}
\item ECOS: Second-order cone program solver~\cite{boyd_scp}, faster for small/medium problems
\item SCS: First-order solver~\cite{boyd_scp}, handles larger problems
\end{itemize}

This baseline validates correctness of our ADMM solver against established SOCP implementations.

\textbf{Method 3: Prior Work.} Implementation of Jia & Evans~\cite{jia_evans} using manifold regression constraints. This provides comparison to recent constrained rotation smoothing work.

\subsection{Parameter Sensitivity}

We analyze sensitivity to key parameters to understand robustness and optimal configurations.

\textbf{Smoothing Weights ($\lambda, \mu$):} We evaluate the impact of first-order vs second-order smoothing:
\begin{itemize}
\item \textbf{Low $\mu$} (0.05): Emphasizes first-order smoothness, good for slow motion
\item \textbf{Balanced} ($\lambda=1.0, \mu=0.2$): Default, balanced approach
\item \textbf{High $\mu$} (0.5): Emphasizes second-order smoothness, good for reducing jerk
\end{itemize}

\textbf{Trust-Region Parameters:} We test different initial $\Delta$ and adaptation strategies:
\begin{itemize}
\item \textbf{Fixed:} $\Delta = 0.2$ rad constant
\item \textbf{Adaptive:} Dynamically adjusted based on $\rho_k$
\item \textbf{Starting values:} $\Delta \in \{0.1, 0.2, 0.5\}$ rad
\end{itemize}

\textbf{ADMM Parameters:} We study convergence speed vs parameter choices:
\begin{itemize}
\item \textbf{Penalty $\rho$:} $\rho \in \{0.5, 1.0, 2.0\}$
\item \textbf{Tolerance:} $\text{tol} \in \{10^{-5}, 10^{-4}, 10^{-3}\}$
\item \textbf{Max iterations:} 500, 1000, 2000
\end{itemize}

\subsection{Theoretical Enhancement Validation}

We empirically validate the theoretical enhancements from Section~\ref{sec:theory}.

\textbf{Warm-starting Evaluation:} We compare convergence speed with and without warm-starting on sequences of varying smoothness:
\begin{itemize}
\item \textbf{High smoothness:} Consecutive problems are similar, warm-starting most effective
\item \textbf{Variable motion:} Moderate benefit from warm-starting
\item \textbf{Abrupt changes:} Limited benefit, but warm-starting does not hurt
\end{itemize}

We measure reduction in outer iterations and total runtime.

\textbf{Adaptive Trust-Region Evaluation:} We compare fixed vs adaptive parameters:
\begin{itemize}
\item \textbf{Consistent conditions:} Similar performance
\item \textbf{Variable acceptance:} Adaptive achieves better stability and faster convergence
\item \textbf{High noise:} Adaptive more robust to poor linearization
\end{itemize}

\textbf{Convergence Rate Estimation:} We validate that estimated convergence rates from $\log \|\delta^k\|$ match theoretical predictions of $\alpha \in [0.3, 0.5]$. We also evaluate using estimated rates for early stopping when sufficient convergence is detected.

\textbf{KKT Condition Monitoring:} We track KKT residual metrics throughout optimization to verify solution quality. Results show:
\begin{itemize}
\item Primal stationarity residual below $10^{-4}$ at convergence
\item Dual feasibility violations negligible ($< 10^{-6}$)
\item Complementary slackness satisfied within numerical tolerance
\end{itemize}

\subsection{Reproducibility}

All experiments are designed to be fully reproducible:
\begin{itemize}
\item \textbf{Code availability:} Implementation released under open-source license
\item \textbf{Random seeds:} Fixed seeds for all experiments (42, 0, 1, etc.)
\item \textbf{Data generation:} Synthetic generation scripts provided
\item \textbf{Parameter logging:} All algorithm parameters logged and reported
\item \textbf{Environment:} Python 3.11, NumPy 1.26, SciPy 1.13, CVXPY 1.4
\end{itemize}

Benchmarking scripts automate running experiments across multiple parameter configurations and generating result tables suitable for publication.
